\documentclass{subfiles}
\begin{document}
    That idea behind asymptotic notation is that of analyzing algorithms,
        with respect to a given resource, more often then not being time,
        without caring about constants.

    We recall the existence of three major notation: the \emph{omega} notation 
        (\(\Omega\)), the \emph{theta} notation (\(\Theta\)) and the big-o notation 
        (\(\mathcal{O}\)). Let \(g(n) \text{ and } f(n)\) be two functions,
        then we define each notation as follow:
        \[\begin{aligned}
            \omegaOf{f(n)} =& \set{g(n)}[ %
                \exists c > 0, n_{0} \ge 0 \text{ s.t. } %
                    g(n) \ge cf(n), \forall n \ge n_{0}%
            ]\\ 
            %
            \thetaOf{f(n)} =& \set{g(n)}[ %
                \exists c_{1}, c_{2} > 0, n_{0} \ge 0 \text{ s.t. } %
                    c_{1}f(n) \le g(n) \le c_{2}f(n), \forall n \ge n_{0} %
            ]\\ 
            % 
            \bigO{f(n)} =& \set{g(n)}[ %
                \exists c, n_{0} \ge 0 \text{ s.t. } %
                    g(n) \le cf(n), \forall n \ge n_{0} %
            ]
        \end{aligned}\]

    Although we give the definition in set notation, 
        we write \(g(n) = \bigO{f(n)}\) to mean \(g(n) \in \bigO{f(n)}\);
        we do similarly for the other notations. Additionally,
        unless stated otherwise, assume we are interested in the time complexity of the algorithm.
\end{document}
